% Standalone problem document, generated from the adjacent README.md.
% Compile with XeLaTeX. Mathematical content is maintained in Markdown.
\documentclass[11pt,a4paper]{article}
\usepackage[fontsize=10.5pt]{scrextend}
\usepackage[margin=22mm,headheight=15pt,headsep=9mm,footskip=11mm]{geometry}
\usepackage{fontspec}
\setmainfont{texgyrepagella}[Extension=.otf,UprightFont=*-regular,BoldFont=*-bold,ItalicFont=*-italic,BoldItalicFont=*-bolditalic]
\setsansfont{texgyreheros}[Extension=.otf,UprightFont=*-regular,BoldFont=*-bold,ItalicFont=*-italic,BoldItalicFont=*-bolditalic]
\setmonofont{texgyrecursor}[Extension=.otf,UprightFont=*-regular,BoldFont=*-bold,ItalicFont=*-italic,BoldItalicFont=*-bolditalic,Scale=MatchLowercase]
\usepackage{amsmath,amssymb}
\usepackage{unicode-math}
\setmathfont{texgyrepagella-math.otf}
\usepackage{xcolor,microtype,enumitem,needspace,fancyhdr}
\usepackage{bookmark}
\definecolor{ink}{HTML}{17344B}
\definecolor{accent}{HTML}{197D80}
\definecolor{muted}{HTML}{596773}
\hypersetup{colorlinks=true,linkcolor=accent,urlcolor=accent,pdftitle={MF-03 - A
uniform disk bound for wave-kernel Padé
approximants},pdfauthor={Open Problems in Numerical Linear Algebra}}
\urlstyle{same}
\setlength{\parindent}{0pt}
\setlength{\parskip}{5pt plus 1pt minus 1pt}
\linespread{1.04}
\setlength{\emergencystretch}{3em}
\widowpenalty=10000
\clubpenalty=10000
\displaywidowpenalty=10000
\allowdisplaybreaks[1]
\setlist{nosep,leftmargin=1.5em}
\providecommand{\tightlist}{\setlength{\itemsep}{0pt}\setlength{\parskip}{0pt}}
\providecommand{\pandocbounded}[1]{#1}
\pagestyle{fancy}
\fancyhf{}
\fancyhead[L]{\sffamily\footnotesize\color{muted}OPEN PROBLEMS IN NUMERICAL LINEAR ALGEBRA}
\fancyhead[R]{\sffamily\bfseries\footnotesize\color{accent}MF-03}
\fancyfoot[L]{\parbox[b]{0.9\textwidth}{\sffamily\fontsize{8}{10}\selectfont\color{muted}Editorial ratings; a bounded search cannot certify that a problem remains unsolved.\\Literature check: 2026-09-11}}
\fancyfoot[R]{\sffamily\footnotesize\color{muted}\thepage}
\renewcommand{\headrulewidth}{0.3pt}
\renewcommand{\headrule}{\hbox to\headwidth{\color{accent}\leaders\hrule height \headrulewidth\hfill}}
\makeatletter
\renewcommand{\section}{\Needspace{5\baselineskip}\@startsection{section}{1}{0pt}{12pt plus 2pt minus 2pt}{4pt}{\sffamily\bfseries\large\color{ink}}}
\renewcommand{\subsection}{\Needspace{4\baselineskip}\@startsection{subsection}{2}{0pt}{9pt plus 1pt minus 1pt}{3pt}{\sffamily\bfseries\normalsize\color{ink}}}
\makeatother
\setcounter{secnumdepth}{0}
\begin{document}
{\sffamily\small\color{muted}Matrix functions and stability\par}
\vspace{3pt}
{\raggedright\hyphenpenalty=10000\exhyphenpenalty=10000\sffamily\bfseries\fontsize{21}{24}\selectfont\color{ink}A
uniform disk bound for wave-kernel Padé approximants\par}
\vspace{7pt}
{\sffamily\small\textbf{Difficulty:} hard\quad\textbf{Importance:} interesting
to specialist\par}
{\sffamily\small\bfseries\color{accent}Status: Solved\par}
\vspace{4pt}
{\color{accent}\hrule height 0.8pt}
\vspace{6pt}
\textbf{Rating rationale:} Hard because this is a focused all-orders
Padé inequality within established approximation theory; its immediate
impact is on specialist wave-kernel error analysis.

\subsection{Resolution --- 2026-09-11}\label{resolution-2026-09-11}

\textbf{Author:} Matthew J. Colbrook, Department of Applied Mathematics
and Theoretical Physics, University of Cambridge. \textbf{Independent
Codex-agent review: PASS for the exact target.}

For every integer \(m\ge1\), the normalized diagonal Padé denominator
for \(\cosh\sqrt z\) is nonzero on \(|z|\le3\) and \(|1-r_m(z)|\le2\)
there. The bound is strict for \(m\ge2\) and sharp for \(m=1\) at
\(z=3\). The analytic tail argument and exact finite certificates cover
every order.

The complete target is resolved. Its former difficulty rating is
historical; the original statement, references and dated audits remain
below.

\textbf{Primary reference:}
\href{https://github.com/ajt60gaibb/OpenProblemsInNLA/blob/main/references/colbrook-matrix-functions-2026-09-11/manuscripts/MF-03.pdf}{complete
authored PDF},
\href{https://github.com/ajt60gaibb/OpenProblemsInNLA/blob/main/references/colbrook-matrix-functions-2026-09-11/manuscripts/MF-03.tex}{standalone
TeX}, \textbf{Theorem 1}.
\href{https://github.com/ajt60gaibb/OpenProblemsInNLA/blob/main/references/colbrook-matrix-functions-2026-09-11/verification/reviews/MF-03-review.md}{Independent
proof review} ·
\href{https://github.com/ajt60gaibb/OpenProblemsInNLA/blob/main/references/colbrook-matrix-functions-2026-09-11/README.md}{Authorship
and submission record}. Verification is independent agent review, not
external human peer review or formal certification.

\subsection{Problem statement}\label{problem-statement}

Let

\[
f(z)=\sum_{j=0}^\infty\frac{z^j}{(2j)!}=\cosh\sqrt z,
\]

where the series defines the entire function without a square-root
branch choice. For each integer \(m\ge1\), let \(r_m=P_m/Q_m\) be its
diagonal Padé approximant at zero: \(\deg P_m,\deg Q_m\le m\),
\(Q_m(0)=1\), and \(Q_m(z)f(z)-P_m(z)=O(z^{2m+1})\). Is it true that the
reduced rational function \(r_m\) has no pole in
\(\{z\in\mathbb C:|z|\le3\}\) and

\[
|1-r_m(z)|\le2\qquad (|z|\le3)
\]

for every \(m\)?

\subsection{Reference and status
evidence}\label{reference-and-status-evidence}

Nadukandi and Higham,
\href{https://eprints.maths.manchester.ac.uk/2651/3/manuscript_nadukandi_higham_wkm_2018_08_01.pdf}{Computing
the Wave-Kernel Matrix Functions}, SIAM J. Scientific Computing 40(6)
(2018), \href{https://doi.org/10.1137/18M1170352}{DOI}, §4.2, Conjecture
4.6, manuscript p.~12. Lemma 4.5 establishes the finite range
\(m\le20\); §4.3 explains its role in backward-error analysis. Searches
for the paper title with ``conjecture'' and for ``Conjecture 4.6'' with
``cosh'' and ``Padé'' located no general proof or counterexample. The
status evidence is therefore weaker than a recent paper explicitly
reaffirming the conjecture.

\subsection{Audit --- 2026-09-10}\label{audit-2026-09-10}

Rechecked
\href{https://eprints.maths.manchester.ac.uk/2651/3/manuscript_nadukandi_higham_wkm_2018_08_01.pdf}{Lemma
4.5 and Conjecture 4.6}: the displayed assertion is proved for
\(1\le m\le20\), while arbitrary \(m\) remains conjectural. Paper-title
and conjecture-number searches found no general resolution. The open
portion still rests on historical primary evidence.
\end{document}
